\section{Conclusion}\label{c:conclusion}

\subsection{What was achieved}

The project delivers a complete, working compiler front end for a language with no standard
written form: a corpus of \CorpusTotal\ statements heard in Yaound\'e, transcribed by hand,
cross-checked and attributed across all three members with the attestation enforced in
software; a hand-written scanner and lexer over five documented regular expressions,
resolving \LexiconTotal\ lexicon entries across \OriginCount\ donor languages into
\TerminalCount\ terminal classes with multiword recognition, accent folding and homograph
protection; a \ProductionCount-production context-free grammar authored in natural
left-recursive form and mechanically transformed in \LedgerSteps\ recorded steps into an
\ExecProductionCount-production LL(1) grammar; FIRST and FOLLOW sets and a parsing table
verified conflict-free by an automated test; and a table-driven predictive parser accepting
\CorpusAccepted\ of \CorpusTotal\ collected statements, returning an explicit derivation on
acceptance and an expected-terminal set plus concrete repair suggestions on rejection.

\subsection{What we learned}

The most useful lesson was that the transformations taught in the course are not
bookkeeping. Eliminating left recursion on \nonterm{Seq} \emph{created} a new common prefix
on \term{SEP} that had to be factored out afterwards --- the two techniques interact, and the
order in which they are applied decides whether the result is LL(1). Seeing that happen on
our own grammar, in the ledger the tool recorded, taught more than the worked examples did.

The second lesson concerns honesty in tooling. It was tempting, each time a statement was
rejected, to add a production until it passed. We resisted, and the two remaining rejections
in Section~\ref{c:results} are documented with the specific LL(1) conflict that prevents a
fix. A grammar that accepts everything distinguishes nothing.

The third was methodological: generating every table in this document from the running
system, rather than transcribing results by hand, caught several places where earlier
documentation had drifted away from the code it claimed to describe.

\subsection{Limitations and future work}

Two constructions --- post-object adverbs and post-verbal clitics --- lie outside the LL(1)
grammar, for the reasons given in Section~\ref{c:results}. The lexicon's dictionary tier
(\LexiconDictionary\ of \LexiconTotal\ entries) is reference-derived rather than
corpus-attested, and is labelled as such. Analysis is purely syntactic: the analyzer
recognises structure, it does not interpret meaning.

Splitting \term{PRON} into full pronouns and clitics would remove both documented rejections
at the cost of re-annotating the pronoun entries against fresh evidence. Beyond that, the
natural next step is a semantic layer: the lexicon already carries glosses and donor
languages, so a gloss-composition pass over the parse tree would turn the analyzer from a
recogniser into a translator. A larger corpus would also let the topic matcher be evaluated
for precision and recall rather than merely cited for evidence.

\subsection{Reproducing this report}

No figure or table here was transcribed by hand. From \texttt{docs/report}:

\begin{lstlisting}[language=bash, numbers=none]
python tools/export_tables.py       # regenerate generated/*.tex from the analyzer
python tools/render_diagrams.py     # regenerate figures/*.png from diagrams/*.puml
latexmk -pdf main-compact.tex       # build this document
\end{lstlisting}

Screenshots are captured separately from the front-end repository with
\texttt{node scripts/\allowbreak capture-report-\allowbreak screenshots.mjs}. Source code
accompanying this report comprises the lexical analyzer (\texttt{core/text.py},
\texttt{core/lexer.py}), the parser (\texttt{core/ll1.py}, \texttt{core/parser.py}) and 158
tests in \texttt{tests/} drawn from the collected data.
